\section{Diseño Algorítmico y Estructuras de Colas de Prioridad}
\label{sec:diseno_algoritmico}

\subsection{Algoritmo de Dijkstra con Cola de Prioridad Abstracta}
El algoritmo de Dijkstra calcula las rutas mínimas sobre grafos con pesos no negativos operando sobre una cola de prioridad genérica $Q$. El procedimiento formal estructurado se describe en el Algoritmo \ref{alg:dijkstra_generico}.

\begin{algorithm}[H]
\caption{Algoritmo de Dijkstra con Cola de Prioridad Parametrizada}
\label{alg:dijkstra_generico}
\begin{algorithmic}[1]
\REQUIRE Grafo dirigido $G = (V, E, W)$, Vértice origen $s \in V$, Estructura de cola $Q$
\ENSURE Distancias mínimas $d: V \to \mathbb{R}^+ \cup \{\infty\}$, Predecesores $\pi: V \to V \cup \{\text{NIL}\}$
\STATE \textbf{Inicialización:}
\FORALL{$v \in V$}
    \STATE $d[v] \leftarrow \infty$
    \STATE $\pi[v] \leftarrow \text{NIL}$
\ENDFOR
\STATE $d[s] \leftarrow 0.0$
\STATE $Q.\text{BUILD}(V, d)$ \COMMENT{O inserción iterativa de cada $v \in V$ con clave $d[v]$}
\STATE $S \leftarrow \emptyset$
\WHILE{\textbf{not} $Q.\text{IS\_EMPTY}()$}
    \STATE $(u, d_u) \leftarrow Q.\text{EXTRACT\_MIN}()$
    \IF{$d_u = \infty$}
        \STATE \textbf{break} \COMMENT{Vértices restantes inalcanzables desde $s$}
    \ENDIF
    \STATE $S \leftarrow S \cup \{u\}$
    \FORALL{$(u, v) \in E$}
        \IF{$v \notin S$}
            \STATE $\text{candidate\_dist} \leftarrow d_u + W(u, v)$
            \IF{$\text{candidate\_dist} < d[v]$}
                \STATE $d[v] \leftarrow \text{candidate\_dist}$
                \STATE $\pi[v] \leftarrow u$
                \STATE $Q.\text{DECREASE\_KEY}(v, d[v])$ \COMMENT{Actualización de prioridad en $Q$}
            \ENDIF
        \ENDIF
    \ENDFOR
\ENDWHILE
\RETURN $(d, \pi)$
\end{algorithmic}
\end{algorithm}

\subsection{Invariantes de Bucle y Prueba de Corrección}
La corrección del algoritmo descansa en el siguiente invariante formal de bucle:

\paragraph*{Invariante:} Al inicio de cada iteración del bucle \texttt{while} (línea 9), para todo vértice $u \in S$, el valor $d[u]$ es exactamente la distancia mínima desde el origen $s$ a $u$, esto es, $d[u] = \delta(s, u)$. Además, para cada vértice tentativo $v \in Q = V \setminus S$, el valor $d[v]$ representa la longitud del camino más corto desde $s$ hasta $v$ cuyos vértices intermedios pertenecen exclusivamente al conjunto de permanentes $S$.

\paragraph*{Demostración por Inducción Matemática:}
\begin{itemize}
    \item \textbf{Base ($|S|=0$):} Antes de la primera iteración, $S = \emptyset$. Para el origen $s$, $d[s] = 0 = \delta(s, s)$ puesto que los pesos son no negativos ($w(e) \ge 0$). Para todo $v \ne s$, $d[v] = \infty$, lo cual satisface la cota superior vacía.
    \item \textbf{Paso Inductivo:} Supóngase que el invariante se cumple para un conjunto $S$ de tamaño $k$. El algoritmo selecciona $u \in Q$ tal que $d[u] = \min_{z \in Q} d[z]$. Si existiese un camino más corto $P$ desde $s$ a $u$, dicho camino debe abandonar $S$ por primera vez a través de una arista $(x, y)$ con $x \in S$ y $y \in Q$. La longitud de este camino sería:
    \begin{equation}
        w(P) \ge \delta(s, x) + w(x, y) = d[y] \ge d[u]
    \end{equation}
    dado que los pesos son no negativos y $u$ minimiza la distancia tentativa en $Q$. Por lo tanto, no puede existir un camino más corto que $d[u]$, implicando $d[u] = \delta(s, u)$.
    \item \textbf{Terminación:} Cuando $Q$ queda vacía o los nodos restantes tienen distancia $\infty$, $S = V$ (para la componente conexa alcanzable), garantizando que $d[v] = \delta(s, v)$ para todo $v \in V$.
\end{itemize}

\subsection{Estructuras de Datos de Cola de Prioridad Evaluadas}

\subsubsection{Min-Heap Binario Indexado (\texttt{IndexedMinHeap})}
En un montículo binario estándar implementado sobre un arreglo contiguo, la operación \texttt{Decrease-Key}$(v, \text{nueva\_clave})$ requiere localizar previamente el nodo $v$, lo cual demanda una búsqueda lineal $O(|V|)$ inaceptable para grafos grandes.

Nuestra implementación desacoplada resuelve este cuello de botella mediante una tabla hash inversa de posiciones $\text{pos}: V \to \mathbb{Z}^+$. Cuando dos elementos en las posiciones $i$ y $j$ del arreglo se intercambian durante las rutinas de flotado (\texttt{\_sift\_up}) o hundido (\texttt{\_sift\_down}), sus entradas en $\text{pos}$ se actualizan de forma atómica en $O(1)$:
\begin{equation}
    \text{heap}[i], \text{heap}[j] \leftarrow \text{heap}[j], \text{heap}[i]; \quad \text{pos}[\text{heap}[i].\text{id}] \leftarrow i; \quad \text{pos}[\text{heap}[j].\text{id}] \leftarrow j
\end{equation}
Con este mecanismo, \texttt{Decrease-Key} ubica el índice en $O(1)$ y ejecuta a lo sumo $\lfloor \log_2 |V| \rfloor$ intercambios hacia la raíz, logrando una complejidad estricta de $O(\log |V|)$.

\subsubsection{4-ary Heap Contiguo (\texttt{DAryHeap})}
El $d$-ary Heap generaliza el montículo binario permitiendo que cada nodo posea $d$ hijos. Para un índice de nodo $i$ (indexación base cero):
\begin{equation}
    \text{padre}(i) = \left\lfloor \frac{i - 1}{d} \right\rfloor, \quad \text{hijo}_k(i) = d \cdot i + k \quad (\text{con } 1 \le k \le d)
\end{equation}
Al fijar $d=4$, la altura máxima del árbol se reduce a $\lceil \log_4 |V| \rceil = \frac{1}{2} \lceil \log_2 |V| \rceil$. Aunque la operación \texttt{\_sift\_down} debe inspeccionar 4 hijos para hallar el menor (requiriendo 3 comparaciones por nivel), la operación crítica \texttt{Decrease-Key} ejecuta el flotado ascendente \texttt{\_sift\_up} con una única comparación por nivel:
\begin{equation}
    T_{\text{Decrease-Key}} = O(\log_4 |V|) = O\left(\frac{\log_2 |V|}{2}\right)
\end{equation}
Adicionalmente, los 4 hijos de cualquier nodo residen en posiciones contiguas del arreglo, lo que maximiza la probabilidad de que se alojen dentro de la misma línea de memoria caché L1/L2 (típicamente de 64 bytes), reduciendo drásticamente los fallos de caché frente a árboles dinámicos con punteros dispersos \parencite{lamarca1996cache}.

\subsubsection{Fibonacci Heap (\texttt{FibonacciHeap})}
El Fibonacci Heap \parencite{fredman1987fibonacci} es una colección de árboles de orden min-heap enraizados en una lista circular doblemente enlazada. A diferencia de los montículos basados en arreglos, no impone una forma rígida de árbol completo. Sus propiedades matemáticas distintivas son:
\begin{enumerate}
    \item \textbf{Inserción Perezosa ($O(1)$ peor caso):} Los nuevos nodos se incorporan a la lista de raíces sin consolidar.
    \item \textbf{Disminución de Clave ($O(1)$ amortizado):} Si al decrecer la clave de $x$ se viola la propiedad de montículo frente a su padre $y = p[x]$, el subárbol enraizado en $x$ se corta inmediatamente y se transfiere a la lista de raíces. Si el padre $y$ ya había perdido previamente un hijo ($\text{mark}[y] = \text{true}$), se desencadena un corte en cascada (\textit{cascading cut}) recursivo hacia los ancestros.
    \item \textbf{Extracción del Mínimo ($O(\log |V|)$ amortizado):} Se extrae la raíz mínima y todos sus hijos se promueven a la lista raíz. A continuación, se ejecuta una rutina de consolidación vinculando raíces de igual grado mediante un arreglo auxiliar de punteros de tamaño $O(\log |V|)$.
\end{enumerate}

\subsubsection{Pairing Heap (Perspectiva Teórica)}
El Pairing Heap es un árbol autoajustable multi-camino que simplifica radicalmente la estructura del Fibonacci Heap al no requerir el almacenamiento de grados ni banderas de corte, logrando tiempos empíricos comparables y $O(\log |V|)$ amortizado para \texttt{Decrease-Key}.

\subsection{Matriz Comparativa de Complejidad Asintótica}
La Tabla \ref{tab:complejidades} resume formalmente las cotas asintóticas temporales y el costo total de ejecución del Algoritmo de Dijkstra según la cola de prioridad empleada.

\begin{table}[H]
    \centering
    \fontsize{10}{12.5}\selectfont
    \renewcommand{\arraystretch}{1.25}
    \setlength{\tabcolsep}{3.5pt}
    \caption{Matriz comparativa de complejidades temporales y espaciales de colas de prioridad en Dijkstra}
    \label{tab:complejidades}
    \begin{tabularx}{\textwidth}{@{}>{\bfseries}l ccccc >{\raggedright\arraybackslash}X@{}}
        \toprule
        \textbf{Estructura} & \textbf{Insert} & \textbf{Extract-Min} & \textbf{Dec-Key} & \textbf{Espacio} & \textbf{Costo Total Dijkstra} & \textbf{Régimen Óptimo en Red Vial} \\
        \midrule
        Arreglo No Ord. & $O(1)$ & $O(|V|)$ & $O(1)$ & $O(|V|)$ & $O(|V|^2)$ & Grafos densos ($|E| \approx |V|^2$). \\
        Binary Min-Heap & $O(\log |V|)$ & $O(\log |V|)$ & $O(\log |V|)$ & $O(|V|)$ & $O((|V|+|E|)\log |V|)$ & Grafos dispersos estándar. \\
        Montículo 4-ario & $O(\log_4 |V|)$ & $O(4\log_4 |V|)$ & $O(\log_4 |V|)$ & $O(|V|)$ & $O((4|V|+|E|)\log_4 |V|)$ & Alta localidad en línea de caché. \\
        Fibonacci Heap & $O(1)$ & $O(\log |V|)^*$ & $O(1)^*$ & $O(|V|)$ & $O(|E| + |V|\log |V|)$ & Óptimo teórico ($|E| \gg |V|$). \\
        Pairing Heap & $O(1)$ & $O(\log |V|)^*$ & $o(\log |V|)^*$ & $O(|V|)$ & $O(|E| + |V|\log |V|)$ & Menor sobrecarga de enlaces. \\
        \bottomrule
    \end{tabularx}
    \vspace{0.15cm}
    \raggedright \fontsize{9}{11}\selectfont \textit{Nota.} El asterisco ($^*$) denota costo amortizado según el método del potencial. Para montículos $d$-arios ($d=4$), la cota es $O((d \cdot |V| + |E|)\log_d |V|)$.
\end{table}
